\section{Introduction}
\label{sec:introduction}
A text editor keeps the document it is editing in memory, and every key
press changes it. The simplest representation, a single contiguous array of
characters, makes reading fast but makes editing slow: inserting one
character near the start of a large file moves every character after it.

Editors have answered this problem in several ways. A gap buffer keeps an
empty region at the cursor, so that typing at the cursor is cheap while
jumping elsewhere is not~\cite{gap}. A piece table records edits as a list
of references into the original file and an append-only buffer. A rope,
the subject of this paper, stores the text in the leaves of a balanced
binary tree~\cite{boehm}.

Section~\ref{sec:design} describes how a rope represents text and how the
basic operations work. Section~\ref{sec:evaluation} compares ropes with the
alternatives on a few common editing patterns.
